% Revision 4 -- proposed abstract, contribution list and conclusion. Every number is from
% paper_rev4/generated/ (simulation, 30 seeds) unless it says emulation.

% ---- Abstract ----------------------------------------------------------------------------------
\begin{abstract}
Link-flooding attacks congest selected links with low-rate, legitimate-looking flows, so an SDN
controller cannot simply filter them; it can only reroute. We formulate attack-aware rerouting as an
online, unsplittable routing problem that minimises the total utilisation in excess of a threshold
over any number of attackers, attacked links and traffic demands, and we derive ResilientSDN as a
myopic drift-plus-penalty approximation of it: a trend detector, a capacity filter that provably
creates no new congested link, a multi-objective path score and a hysteresis deadband. We evaluate it
in Mininet on a fat-tree and three real backbones (NSFNET, Internet2, G\'EANT) and, with the same
decision code, in a flow-level simulator over 30 seeds, eleven attack patterns plus multi-link and
adaptive (rolling) attacks, UDP and TCP traffic, and baselines including ECMP, $k$-path multipath and
two traffic-engineering LPs. ResilientSDN cut the time an attacked link spends above the threshold from
121\,s to 35\,s per attack on average, with 2.3 flow moves per run against 388 for threshold-reactive
rerouting, and kept victims' TCP flow-completion times within 2.0--3.3\,s (static routing:
7.9--17\,s). An improved variant that scores paths by their bottleneck instead of summed utilisation
reduced the time to 15\,s and removed a failure on sparse WANs. Fractional LP-based traffic
engineering achieved less congestion (11--24\,s) at the price of splitting flows. To test whether
these results generalise beyond hand-made scenarios, we also drew 1{,}945 random instances of the
problem (real and random topologies, gravity-model demands, one to three attacked links, four attack
behaviours) and compared every system with the instance optimum of a multicommodity LP. Over them,
the variant with a duty-cycle trigger left a median of 40\,\% of the avoidable excess (ResilientSDN:
70\,\%, the TE LP: 9\,\%) and closed most of the trend detector's pulse loophole; routing-aware rolling
attackers remain the case in which rerouting can do worse than not reacting.
\end{abstract}

% ---- Contributions (introduction) -----------------------------------------------------------
\begin{itemize}
\item A formulation of attack-aware rerouting for multiple attackers, multiple and adaptive target
  links, arbitrary demands and topologies (Sec.~\ref{sec:formulation}), its NP-hardness, two LP
  relaxations, an exact per-instance optimum $\Phi^\star$ that every controller is measured against,
  a cut condition that tells when no controller can win (Proposition~\ref{prop:cut}), and three
  guarantees for ResilientSDN: its capacity-filtered greedy step never creates
  secondary congestion (Proposition~\ref{prop:monotone}), its cooldown bounds every flow's moves
  whatever the attacker does (Proposition~\ref{prop:churn}), and an attacker that stays undetected can
  keep a link congested for at most a known fraction of the time (Proposition~\ref{prop:evasion}).
\item A randomised evaluation over the scenario space (Sec.~\ref{sec:sweep}): 1{,}945 instances on
  eight real and two families of random topologies with random demands, attack intensities, numbers
  of targets and attack behaviours, which tests the conclusions of the hand-made scenarios and shows
  where they do not hold.
\item ResilientSDN, an online drift-plus-penalty approximation of that problem, whose decision logic
  (Algorithm~\ref{alg:rsdn}) is a single module shared by the Ryu controller and the simulator.
\item An evaluation in Mininet (four topologies, eleven attack patterns) and in a validated flow-level
  simulator (29\,280 runs, 30 seeds, multi-link and adaptive attacks, TCP), with paired,
  multiplicity-corrected statistics and multipath and LP traffic-engineering baselines.
\item An honest account of where it fails --- pulse attacks, adaptive attackers, sum-drift holds on
  sparse WANs --- and of which extension fixes what (ResilientSDN-v2/v3 ablation).
\end{itemize}

% ---- Conclusion --------------------------------------------------------------------------------
\section{Conclusion}
ResilientSDN reroutes around link-flooding attacks with few, verified moves: across four topologies
and fourteen scenarios an attacked link stayed above the threshold for 35\,s per attack on average,
against 121\,s with static routing and 56\,s with threshold-reactive rerouting, without the latter's
oscillation, and it kept TCP victims' completion times close to ideal. The generalised formulation explains both its guarantees (no secondary congestion
from filtered moves) and its failures (the sum-of-utilisation drift that holds flows on sparse WANs,
fixed by bottleneck drift). Two results bound the claim: fractional traffic engineering with the same
measurements does better, at the cost of splitting flows and more rule changes, and attackers that pulse
or re-target adaptively still keep links congested for a large part of an attack. A randomised sweep
over 1{,}945 instances of the general problem confirmed most of the ranking of the hand-made scenarios
--- except that static $k$-path splitting comes out ahead of ResilientSDN in the median, while being
worse than static routing in a third of the instances --- and the pulse loophole that
Proposition~\ref{prop:evasion} predicts; it showed that a duty-cycle trigger closes
most of it, and showed that against a routing-aware rolling attacker every reactive system,
ResilientSDN included, is sometimes worse than not rerouting at all. In the scenario simulations the
duty-cycle trigger cut the pulse attack S7 from 32 to 6\,s above the threshold without a single move
in the absence of an attack, and brought the combined variant (v3) to 11.8\,s per attack. Future work
is to evaluate the v2/v3 extensions in emulation, to detect re-targeting so that rerouting stops when
the attack follows the victims, to add the multi-link, adaptive and TCP scenarios to the Mininet
testbed, and to combine single-path rerouting with limited splitting where it pays off.
